\section{Discussion}\label{sec:discussion}

\subsection{The answer}

On holiday, do Spaniards go back to the villages their families left? The data cannot follow a person to their own village, but they answer the question at the level of cities and zones, and the answer there is yes in a specific sense. On every night tested, cities send more of their residents to the provinces where more of them were born. In August this link becomes clearly stronger for two groups, the residents who arrived before 1981 and those born in villages, and, against February, the increase is found only in stays inland; whether the average link grows depends on the baseline. At the same time, the most emptied zones, which send people away on an ordinary night, receive more than they send on almost every holiday night. The beaches take a large part of the summer and shape the network as a whole, but the inland stays still follow the map of the exodus.

\subsection{What was expected and what was not}

I expected the summer to be about the coast, and in one sense it is: the network splits into cities that send and towns that receive, its largest receivers are coastal, and most of the change follows from the new degrees. I did not expect the roots to be so visible on ordinary nights, nor the average holiday increase to be so small. The result that surprised me most is that the average hides two opposite movements. The exodus generation pulls harder in August, while recent migrants pull less against February, although that second result is weaker against other baselines.

The two holidays also bring people back in different ways. In August the emptied zones become receivers mainly because fewer of their residents leave, and the stock of non-primary dwellings matters for the summer rise in visitors but hardly at Christmas. At Christmas the net gain per resident rises as well. A single holiday average would have mixed the two.

\subsection{What is new}

As far as I could find, no earlier study compares a national holiday network with where its residents were born while also testing its structure against null models. Doing both showed two things: most of the structural change on holiday is forced by the new degrees, and the link to birthplaces only shows clearly once residents are split by when they arrived and where they were born. Earlier work compared holiday travel with migration between 31 provinces without controls for distance \citep{li2016spring}, or measured return tourism in one Spanish region \citep{garciagonzalez2024retorno}. There is also a practical side: on holiday nights the zones that lost most people receive more people than they send, and these short stays do not appear in the population register.

\subsection{Limitations}

\begin{itemize}
\item \textbf{The data are estimates.} The overnight counts come from one mobile-phone operator, scaled up to the population, and each phone's home zone is inferred from where it usually sleeps. People who spend long summers in the village may be counted as residents there.
\item \textbf{The test is about cities, not people.} The roots network gives the birthplaces of a city's residents, not of the people who travel. A city with more residents born in a province sends more people there, but I cannot tell whether the travellers are those residents, their children or their neighbours. The census records the year of arrival in the province, not in the city, and the birthplace by province, not by village.
\item \textbf{Zones group small municipalities.} More than a third of the 2,435 zones (908) group several municipalities, so a ``village'' is often a group of villages. The village results hold in zones made of a single municipality, but grouping still hides what happens in each village.
\item \textbf{Six months of one year.} The data cover February, March, August, October, November and December 2024, with checks on 2022 and 2023. Single nights are compared, not whole seasons, and causes are not identified.
\item \textbf{Some sizes depend on choices.} The sign of the generation effects is stable, but their size depends on how city--province cells with no residents are treated (Table~\ref{tab:zero}).
\end{itemize}

\subsection{Lessons for future work}

Looking back, three choices made the most difference. First, single nights: an average over a holiday period mixes the people who leave with those who come back. Second, degree-preserving null models before any interpretation: the denser holiday networks changed almost every property, and most of those changes turned out to follow from the degrees. The rich club needs particular care: its raw value is the result, and a ratio to rewired networks close to one does not mean there is no club \citep{zhou2007structural}. Third, splitting the roots by group, because the average hid groups that moved in opposite directions. If I continued, the next step would be to replace provinces of birth by municipalities of birth, which would need linked or survey data at village level, and to follow the same nights over several years to see whether the pull of the exodus fades as its generation ages.
